\documentclass[table]{exam}

% Preamble
\pagenumbering{gobble}
\pagestyle{empty}

\usepackage[utf8]{inputenc}
\usepackage{amsmath}

\title{A stumble through Pico's mechanical houses}
\author{JellyWellyBelly}
\date{}

\begin{document}

\maketitle

% History on how hard work rewards you in the long run

\section*{What is all of this about?}

Welcome to the story of Edgar, an adventurous, young man in his random walk across the magic town of Pico, a small whimsical village with abnormal inhabitants and even stranger customs. As we dress ourselves with Edgar's clothes, put on he's shoes and learn how every gear, sprocket, machine and even entire systems work on the undergrounds of this village, we can slowly unveil the '\textit{How to these work?}' feeling behind the houses.

\section*{On a more personal note...}

This document serves as a repository of programming problems for anyone to learn, practice or have fun with at their heart's content. If it helped a single person throughout it's existence, then it has accomplished it's goal. This contains exercises that I picked online, in books or that I made up across my exploration during my career and all attributions will be made when required. This may be distributed under all conditions as, education, in my mind, should be \textbf{ALWAYS} free, easily accessible and understood my the most amount of people possible.

The solutions for all exercises should be found in the repository where this document is publicly available. 

Sometimes, it may happen that there are several solutions in one file and the person reading it may have to change the code to see the different solutions in action. I strive to have the best code I can produce at my current expertise and usually in the languages C and Python, with an emphasis on the former. The questions at hand and their instructions are prepared for C only and should be adapted for any other language (in terms of their constructs) when needed.

I decided to write a little short story around it just to make it fun to write and to challenge myself to do something different, since I plan on expanding this occasionally. I'm a below average writer and I will construct a more bare bones, maybe interactive copy of this.

\newpage

\begin{center}
\vspace*{\fill}
\section*{Wadding out of the confort zone}
\vspace*{\fill}
\end{center}

\newpage

\section*{Bare bones exercises - No whimsy :(}

Exercises are written in a roughly increasing order of difficulty wiithin each section. No input validation in needed unless specified. Also, don't take into consideration security issues, as the point of these exercises is learning programming, specifically.


\subsection{I/O}

All of these exercises are possible using "printf" and "scanf" functions only. It is possible to solve these without the use of a single variable, 

\begin{enumerate}

\item Write a program that outputs the string "Hello World!".

\item Write a program that asks a user their name and outputs a greeting with it's name. The interaction should be as follows:

\begin{verbatim}
Write your name: Jelly
Hello Jelly!
\end{verbatim}

The user above only needs to write "\textit{Jelly}", the remaining interaction should be done by your program.

\item Write a program that receives a real number and outputs it with only 1 decimal place without rounding. Here are some examples of some inputs and their corresponding output.

\begin{verbatim}
In: 3.1415
Out: 3.1

In: 3
Out: 3.0

In: -2.71
Out: -2.7

In: 5.0
Out: 5.0
\end{verbatim}

\item Write a program that receives an ASCII character and outputs it corresponding ASCII value. Assume that the input is correct. Please follow the examples below:

\begin{verbatim}
In: A
Out: 65

In: }
Out: 125
\end{verbatim}

\item Now perform the inverse function of the one implemented above.

\begin{verbatim}
In: 65
Out: A

In: 125
Out: }
\end{verbatim}

\item As a last exercise, do a conversion between a user supplied positive integer and it's corresponding hex value, following the examples below:

\begin{verbatim}
In: 0
Out: 0x0

In: 3
Out: 0x3

In: 28
Out: 0x1c

In: 2026
Out: 0x7ea
\end{verbatim}

\subsection{If's}

Now, with the power of the mighty IF construct 

\item Read a positive

\end{document}
